\section{Related work}

The closest econometric comparison is \citet[Proposition 2.3, equation (2.4), Theorem B.8(ii), Lemma B.9, and Theorem 4.9]{Cytrynbaum2026Limits}. His excess-variance identity expresses the design-dependent part of Horvitz--Thompson variance as scaled expected squared prognostic imbalance, the criterion used in \cref{obj:def:criterion}. Writing \leanref{sym:n}{\ensuremath{n}} for sample size, \leanref{sym:d}{\ensuremath{d}} for covariate dimension, and \leanref{sym:s}{\ensuremath{s}} for smoothness, his correctly specified pairwise model admits a structured Gram--Schmidt Walk (GSW) upper bound of order \((d^2\log n/n)^{s\wedge1}\), supported by the pooled Mat\'ern approximation calculation. His lower bound \(n^{-2s/d}\) concerns a full-dimensional smooth class; its bivariate specialization has exponent \(s\). For independent uniform covariates and the centered order-two pooled Sobolev class in \cref{obj:def:sobolev-class}, the present characterization matches upper and lower bounds across sample size and dimension at order \(\min\{1,(B/n)^s\}\), where \leanref{sym:B}{\ensuremath{B=\binom{d}{2}}} is the number of coordinate pairs, throughout the \leanref{S-1}{parameter range \(n,d\geq2\) and \(0<s\leq1\)}. The comparison therefore concerns the logarithmic factor and the explicit interaction-count dependence for this specified class; \cref{obj:prop:published-class-inclusion} supplies the class-transfer relation.

A precise comparison with random-matrix theory identifies an existing route to saturation. Applied to independent, centered, unit-variance cosine coordinates under uniform sampling, \citet[Theorem 1.2(2a)]{KoganNandyHuang2024Extremal} implies the following spectral limit. Let \(G\) denote the sample Gram matrix of the pair-product features using one cosine frequency per coordinate, normalized by \(B\). As \(n,d\) tend to infinity with \(n/B\to\rho\in(0,1)\), where \(\rho\) is the limiting sample-to-feature ratio, its smallest eigenvalue converges almost surely to \((1-\sqrt{\rho})^2\). This positive spectral edge yields a constant-order signing lower bound for the normalized one-frequency submodel. The additional scope of \cref{obj:thm:capacity-answer,obj:thm:log-free-bivariate-capacity} is the matched finite-sample characterization, including the frequency-rich regime \(B/n\to0\), together with attainment by the same assignment procedure for every smoothness in the stated range. The independent coefficient prior in \cref{obj:def:cosine-prior} supplies the growing frequency family for that subcritical comparison.

The criterion also belongs to the minimax experimental-design tradition. \citet[Sections 2.2--2.4]{Kallus2018} connects restrictions on conditional outcome functions to worst-case imbalance and optimal randomized designs, while \citet[Sections 3--5]{Kallus2020} studies mixed assignment distributions and randomization constraints. Here, \cref{obj:def:design-class,obj:def:criterion} specifies fair, outcome-independent assignment laws and places the worst-function comparison outside the expectation over covariates and assignment. The lower bound in \cref{obj:thm:full-design-lower} applies to this entire admissible design class. Its role is to characterize the best achievable excess variance for the fixed Horvitz--Thompson criterion under the pooled smoothness restriction.

The attaining procedure connects this statistical objective to randomized balancing. The GSW design of \citet[Lemma 6.1 and Theorem 6.3]{HarshawSavjeSpielmanZhang2024GSW} combines assignment fairness with covariance control for augmented covariate features. The Edge-Walk method of \citet[Section 2.1]{LovettMeka2015EdgeWalk} develops constrained random walks for partial coloring. The procedure in \cref{obj:def:ordered-boundary-rounding} uses finite randomized boundary moves with an ordered family of balancing constraints. Its statistical application in \cref{obj:def:capacity-handle} uses a feature ordering independent of smoothness, allowing one assignment law to attain the comparison scale simultaneously across the stated Sobolev classes.

\citet[Lemma 4.2 and Proposition 5.1]{Chen2026} studies nonlinear feature balance through a normalized feature Gram matrix \(K\). For \(\phi\in(0,1)\), the design has assignment covariance bounded by \(\{\phi I+(1-\phi)K\}^{-1}\), and a fixed-schedule ridge approximation bound controls the actual Horvitz--Thompson variance. That comparison is conditional on the realized schedule and a chosen nonlinear feature map. The present result instead gives a population minimax capacity matched over an explicit pooled Sobolev class, with a single smoothness-independent construction and a converse for every admissible design.

Discrepancy-based integration provides a related use of randomized balancing with distinct sampling budgets and error criteria. \citet[Theorem 1.1 and Section 3.1]{BansalJiang2024QMC} starts from \(n^2\) independent uniform points and partitions them into \(n\) sets of \(n\) points. For each retained set, the mean squared integration error is bounded by \(\widetilde O_d(\sigma_{\mathsf{SO}}(f)^2/n^2)\), where \(f\) is the integrand and \(\sigma_{\mathsf{SO}}(f)\) is its smoothed-out variation; the displayed order suppresses logarithmic factors. \citet[Theorem 1.1, Section 2.2, and Theorem 3.3]{ChenJiangKirk2025HighDimensionalQMC} retains this quadratic input pool and develops bounds for weighted Fourier classes and low superposition dimension, with explicit coordinate-weight and dimension dependence. \citet[Theorem 1.2]{EzeunalaJhaJiang2026QMCII} obtains the smoothed-out-variation mean squared error bound while retaining \(n\) points from \((1+\varepsilon)n\) independent uniform inputs with high probability, for fixed \(\varepsilon\in(0,1)\); its constants and logarithmic factors depend on dimension and \(\varepsilon\).

Other thinning results emphasize different forms of approximation. The online procedure of \citet[Theorems 1--2]{Dwivedi2019} retains at least one point from every two consecutive uniform inputs and controls unnormalized counting discrepancy over axis-parallel rectangles. Kernel thinning in \citet[Theorem 1 and Corollary 1]{Dwivedi2024} compresses an input approximation into a smaller subset and bounds maximum mean discrepancy, equivalently worst-case integration error over the associated reproducing kernel Hilbert-space unit ball, under the stated regularity and input-radius conditions. In the present experiment, all \(n\) original units receive treatment signs, and performance is measured by expected squared prognostic imbalance over a pooled nonperiodic Sobolev class. These sampling, function-class, and loss specifications make the comparison with integration results precise.

Finally, \citet[Theorems 2.2--2.3 and 3.1]{SudijonoDobribanTchetgen2026Minimax} studies joint minimax choice of design and estimator for finite-population treatment-effect estimation with binary or bounded potential outcomes. For binary outcomes, Bernoulli randomization paired with a nonlinear estimator attains their minimax risk, for which they derive a second-order expansion. The present optimization fixes the Horvitz--Thompson estimator and varies covariate-dependent assignment under a smooth prognostic class. It characterizes the excess-variance gains available through assignment for that estimator and model.
